% =====================================================================
% sections/conclusion.tex — shared (long + short)
% =====================================================================

\section{Conclusion}
\label{sec:conclusion}

This paper revisits the long-run effects of Indonesia's Village Midwife
Program on adult human capital using a richer reconstruction of the
analysis sample, four standardized outcome families, and a layered
identification stack. Three empirical findings emerge.

First, the central effect of the program on children exposed in utero
or before age three is small and statistically imprecise on the four
composite indices in the headline TWFE specification, with point
estimates ranging from essentially zero on the cognition index to
\CoefDepTWFE\ standard deviations on the depression index. The
imprecision is real and not an artifact of multiple comparisons:
parallel-trends tests reject for three of four outcomes and
\textcite{anderson2008}-style sensitivity bounds indicate that
even modest selection on unobservables would erase the estimates.

Second, when the within-mother specification is used to sweep out
time-invariant village confounders that drove non-random program
placement, the depression-index point estimate nearly doubles to
\CoefDepSibFE~SD (SE \SEDepSibFE). This is the cleanest piece of
evidence that the program improves long-run mental health. The pattern
is consistent with downward bias in the naïve TWFE estimator from
placement of midwives in poorer or sicker villages, exactly the
concern \textcite{frankenberg2005} raised about earlier studies of the
program. Health-index estimates also rise modestly under sibling fixed
effects, though they remain statistically imprecise.

Third, the gender-specific pattern documented in earlier work survives
the redesign: health gains concentrate among girls. This is not a
mechanical artifact of the new sample or estimator; the result holds
across TWFE, sibling FE, and the
\textcite{callawaysantanna2021} specifications.

Taken together, the picture that emerges is more cautious than the
ones drawn by \textcite{ahsan2021} and the predecessor DEPII version
of this analysis. Indonesia's Village Midwife Program produced
detectable but smaller effects on adult human capital than recent
staggered-DiD work has implied; once the analysis is honest about
non-random placement and pre-trend violations, the most credible
estimates are concentrated in domains and subgroups where prior
mechanism-based predictions are strongest---girls' health, and the
mental-health margin captured by depression symptoms a generation
later.
